\subsection*{Binary Relations}
\label{subsec:binary-relations}

\begin{definitionbox}{Definition (Binary Relation)}
\begin{definition}[Binary Relation]
\label{def:binary-relation}
Let $A$ be a set.
A \textbf{binary relation on $A$} is a rule $R$ that determines, for each
ordered pair $(a,b)$ with $a,b \in A$, whether $a$ and $b$ stand in the
relation.  We write $a \mathrel{R} b$ when the relation holds, and
$\lnot(a \mathrel{R} b)$ when it does not.
\end{definition}
\end{definitionbox}

\begin{remark*}[Standard quantified statement]
\[
R \text{ is a binary relation on } A
\quad\Longleftrightarrow\quad
R\subseteq A\times A.
\]
\end{remark*}

\begin{remark*}[Definition predicate reading]
$R$ is a binary relation on $A$ exactly when every related pair has both
coordinates in $A$.
\end{remark*}

\begin{remark*}[Negated quantified statement]
\[
\neg\operatorname{BinaryRelation}(R,A)
\Longleftrightarrow
R\not\subseteq A\times A.
\]
\end{remark*}

\begin{remark*}[Interpretation]
In set-theoretic language, a binary relation on $A$ is a subset
$R \subseteq A \times A$.  The notation $a \mathrel{R} b$ treats the relation
as a predicate on pairs, which is the algebraic point of view used here.
\end{remark*}

\begin{dependencies}
\begin{itemize}
  \item \hyperref[def:carrier-set]{Carrier Set}
\end{itemize}
\end{dependencies}
